\subsection{Representação Binária}
\label{sec:binario}

A eletrônica digital baseia-se no bit (binary digit), a menor unidade de informação, que pode assumir apenas dois estados discretos: 0 ou 1. Fisicamente, esses estados são traduzidos por níveis de tensão elétrica, onde o "0" (falso) geralmente corresponde ao aterramento ou ausência de sinal, e o "1" (verdadeiro) a uma tensão positiva, como 5V. Toda a complexidade dos computadores modernos é construída sobre essa abstração lógica binária.

\subsubsection{Sistema Numérico Binário}

O sistema numérico binário é um sistema de base 2, no qual qualquer número é representado utilizando apenas dois dígitos: 0 e 1. Diferentemente do sistema decimal, no qual cada posição representa potências de 10, no sistema binário cada posição representa potências de 2.

Assim, um número binário pode ser interpretado como uma soma ponderada de potências de 2. Por exemplo:

\[
1011_2 = (1 \cdot 2^3) + (0 \cdot 2^2) + (1 \cdot 2^1) + (1 \cdot 2^0)
\]

\[
1011_2 = 8 + 0 + 2 + 1 = 11_{10}
\]

Essa representação é fundamental na eletrônica digital, pois permite mapear diretamente estados físicos de circuitos elétricos (ligado/desligado) para valores numéricos.

Além disso, o sistema binário é a base para a construção de operações aritméticas em hardware, como soma, subtração e comparação, que podem ser implementadas exclusivamente com portas lógicas.